\clearpage\thispagestyle{empty}
\pdfbookmark[0]{Quellen zu Komet (10. September 1811)}{comet1811sep10-sources}
\begin{center}\textbf{Quellen zu Komet (10. September 1811)}\end{center}
\noindent Diese KI-Übertragung ins heutige Deutsch folgt dem Brief und dem Zusatz von Carl Friedrich Gauß in \href{https://gdz.sub.uni-goettingen.de/id/PPN236007467}{Gauß’ \textit{Werken}, Bd.~VI} (Göttingen, 1874), S.~333--335. Sie erschienen unter der Überschrift \textit{Beobachtungen und Elemente der Bahn des Kometen von 1811 und Zusatz zu seiner Theoria motus corp. coelest.} im \textit{Astronomischen Jahrbuch für das Jahr 1814}, herausgegeben von Johann Elert Bode (Berlin, 1811). Der Beitrag beginnt auf S.~\href{https://api.digitale-sammlungen.de/iiif/image/v2/bsb10538343_00262/full/1600,/0/default.jpg}{254}, setzt sich auf S.~\href{https://api.digitale-sammlungen.de/iiif/image/v2/bsb10538343_00263/full/1600,/0/default.jpg}{255} und \href{https://api.digitale-sammlungen.de/iiif/image/v2/bsb10538343_00264/full/1600,/0/default.jpg}{256} fort und endet auf S.~\href{https://api.digitale-sammlungen.de/iiif/image/v2/bsb10538343_00265/full/1600,/0/default.jpg}{257}. Das digitalisierte Exemplar stammt aus der Bayerischen Staatsbibliothek.
\par\smallskip\noindent Das Jahrbuch nennt neun Beobachtungen von Zach und eine Vorausberechnung für 3 Monate. Die Werkausgabe nennt neue Beobachtungen und 4 Monate. Das Jahrbuch enthält eine redaktionelle Fußnote zur Zeit der Bedeckung und setzt auf S.~257 $z$ anstelle von $x$. In der Werkausgabe steht $-z$ anstelle von $x$. Das kleine Zeichen links oberhalb der 2 im Zähler der Formel für $\xi$ ist in der Werkausgabe zu sehen und wird als bearbeitbare Konturen wiedergegeben. KI-Übertragung und KI-Satz.
\clearpage
